% Prefix-balance modeling; no new algorithm implementation.
\section{Catalan 与投票问题：把前缀限制变成反射双射}\label{knowledge-ballot-reflection}
\index{Catalan}\index{ballot theorem}\index{reflection principle}
\textbf{先问限制的是最终数量，还是每个前缀。}
把一种操作记为 $+1$，另一种记为 $-1$，高度就是两种操作数量之差。
固定 $p$ 个正步、$q$ 个负步，共有 $\binom{p+q}{q}$ 个次序；这只固定终点，未限制中途高度。
合法括号串要求高度始终非负且最终为零；从空栈开始的入栈/出栈操作序列也有同样条件。
这只计操作类型序列，不自动计带标签元素的任意排列、操作耗时或额外栈约束。

\subsection{弱前缀条件：减去首次走到负一的坏路径}
假设 $p\ge q\ge0$，从高度 0 出发，每个非空前缀高度都要求 $\ge0$。
坏路径必然有第一次到达 $-1$ 的时刻；只把从起点到该时刻的所有正负步交换，后缀不变。
该前缀原来负步比正步多 1，交换后总正步变成 $p+1$、负步变成 $q-1$。
反过来，任意这样的新路径终点为 $p-q+2\ge2$，必然第一次到达 $+1$；
把到此处的前缀再交换，就唯一恢复一条坏路径。
因此这是双射，而不只是“看起来对称”，合法数为
\[
W(p,q)=\binom{p+q}{q}-\binom{p+q}{q-1}.
\]
约定 $\binom Nk=0$ 当 $k<0$ 或 $k>N$。若 $p<q$，终点已非法，答案直接为 0，
不能继续使用可能给负数的差式。$p=q=0$ 的空序列计 1。
反射必须截止于\emph{第一次}越界；翻转整条路径或任意一个越界前缀都不构成上述逆映射。

取 $p=q=n$，得到括号串、Dyck 路径的 Catalan 数
\[
C_n=\binom{2n}{n}-\binom{2n}{n-1}
   =\frac1{n+1}\binom{2n}{n},\qquad C_0=1.
\]
这里 $\binom{2n}{n-1}=\binom{2n}{n+1}$，与前文差分公式相同。
这些经典计数对象可参见 Stanley 的
\href{https://math.mit.edu/~rstan/ec/catalan.pdf}{Catalan 习题 19(h),(i)}；
使用时仍需先证明题目对象与路径一一对应。

\subsection{严格领先：先取走必需的第一步}
若要求每个非空前缀高度 $>0$，且 $p>q$，第一步必须为正步。
删去这一步，并将剩余路径的高度整体减 1，得到 $p-1$ 个正步、$q$ 个负步的弱条件，故
\[
S(p,q)=\binom{p+q-1}{q}-\binom{p+q-1}{q-1}
      =\frac{p-q}{p+q}\binom{p+q}{q}.
\]
当 $p+q>0$ 且 $p\le q$ 时答案为 0；空路径若按“所有非空前缀”定义，则真空成立、单独计 1，
不能把它代入含 $p+q$ 分母的式子。
例如 $p=3,q=2$：弱条件有 5 种，严格条件只有 $+++--$ 和 $++-+-$ 两种。
“过程中可以打平”和“始终领先”不是同一条件，不能只检查最终 $p>q$。

\newpage
\subsection{初始有余额：反射起点，而非只改终点}
从整数高度 $h\ge0$ 出发，走 $L\ge0$ 个 $\pm1$ 步，到达 $e\ge0$，全程不低于 0。
先检查 $|e-h|\le L$ 且 $L+e-h$ 为偶数，否则答案为 0。
令 $k=(L+e-h)/2$ 为正步数，总路径数为 $\binom Lk$。
坏路径第一次到达 $-1$ 之前，按高度变换 $y\mapsto-2-y$ 反射前缀，
把起点 $h$ 变成 $-h-2$，越界点 $-1$ 不动，终点仍为 $e$。
由于新起点在 $-1$ 下方、终点非负，每条新路径都必经 $-1$，在第一次到达处反射回来就是逆映射。
新路径正步数为 $k+h+1$，故
\[
B(L,h,e)=\binom Lk-\binom L{k+h+1}.
\]
例：$L=4,h=2,e=0$，四种无约束次序中只有 $---+$ 越界，答案为 $4-1=3$。
$L=0$ 时只有 $h=e$ 的空路径计 1。若下界为整数 $b$，先把起终点减去 $b$；
若要求整数高度严格大于 $b$，则下界应取 $b+1$。

\textbf{适用边界。} 单个常数下界和单位 $\pm1$ 步使“第一次碰到边界”可逆。
再加上界、禁止位置、加权路径、非单位步长或斜率边界时，不能照搬两项差式；
可能需要状态 DP、重复反射或另一种专门计数。单步能跳过 $-1$ 时，上述首次命中论证已经失效。

\subsection{取模时优先使用差式，除法条件不能省}
差式只要求能正确计算两个组合数，不要求 $n+1$、$p+q$ 可逆。
例如 $C_2=2$，模 6 仍为 2；把 $\binom42$ 先取模得到 0，再除以 3 无法恢复答案，
因为 3 没有模 6 逆元。即使模数是素数，分母也可能被模数整除。

选择见第~\pageref{knowledge-binomial-choice}~页（\ref{knowledge-binomial-choice}~节）：
素数模数且所需最大上标小于模数时可用
\hyperref[compact-Binomial]{Binomial}（第~\pageref{compact-Binomial}~页，\ref{compact-Binomial}~节）；
跨过素数模数时考虑
\hyperref[compact-Lucas]{Lucas}（第~\pageref{compact-Lucas}~页，\ref{compact-Lucas}~节），
合数模数可按资源预算考虑
\hyperref[compact-ExLucas]{ExLucas}（第~\pageref{compact-ExLucas}~页，\ref{compact-ExLucas}~节）。
后两者都有预处理空间/时间限制，不意味着任意大模数都可直接建立表。
先以有符号或更宽类型检查负下标、不可达性及 $p+q,2n,L+e-h,k+h+1$ 的溢出，
再调用接收无符号参数的接口；最后将两个模值之差规范化。

\newpage
